\chapter{Annexe B — Exemples d'usage du MEF dans la littérature}

\section{Objet de l'annexe}
Le modèle d'évaluation fonctionnelle (MEF)~\cite{autissier2008} a été mobilisé dans plusieurs
secteurs pour diagnostiquer une fonction support — en particulier le service systèmes
d'information — lorsque les seuls indicateurs financiers s'avèrent insuffisants. Cette annexe
reprend, sur le mode narratif, des \textbf{illustrations} déjà publiées dans les ouvrages et études
de cas commentés par Autissier et Delaye, et reprises dans les travaux sur Naftal~\cite{meghari2014}.
Elle ne rapporte pas de nouvelles enquêtes : elle montre comment un outil comme GIA peut s'inscrire
dans la même logique d'\textbf{amélioration continue} que ces diagnostics.

\section{Cas d'une entreprise de BTP et d'aménagement d'espaces verts}
Une PME française du BTP spécialisée dans des aménagements complexes (terrains de golf, espaces
verts) souffrait de la circulation de \textbf{documents papier} sur chantier, avec détérioration et
perte d'information. Le dirigeant souhaitait un outil de gestion de chantiers couvrant conception,
affectation des équipes, suivi des dépenses et de l'avancement. Pour «\,mettre sous tension\,» le
service SI et révéler les lacunes en gestion de projet, le dirigeant a fait appel au MEF.

Les résultats du diagnostic ont conduit à proposer un \textbf{parcours de formation} à la gestion
de projet et à la rédaction de cahiers des charges en optique MOA, en impliquant plusieurs types
d'utilisateurs du futur système. Ce cas illustre deux enseignements transposables à Naftal : (1) la
\textbf{formalisation} précède souvent l'outil ; (2) l'évaluation d'une fonction support peut
déboucher sur des actions de \textbf{compétences} et de \textbf{processus}, pas seulement sur un
achat logiciel.

\section{Cas d'une entreprise industrielle et ratio SI / chiffre d'affaires}
Dans une entreprise industrielle d'envergure moyenne, le contrôleur de gestion a calculé le ratio
entre les dépenses de SI et le chiffre d'affaires (de l'ordre de quelques pour cent). Sans
référentiel, le directeur ne pouvait juger si ce niveau était «\,bon\,» ou «\,élevé\,». Le MEF a été
appliqué avec un accent sur le \textbf{support structurel} (fonctionnement, enjeux, ressources).

Le diagnostic a mis en évidence un \textbf{effectif SI} supérieur aux benchmarks utilisés dans
l'étude, ce qui a ouvert une réflexion sur la \textbf{justification} des ressources au regard de
l'activité réelle et des attentes des utilisateurs. L'enseignement pour une direction comme la DCSI
est double : d'une part, les KPIs «\,métier\,» (tickets, délais, satisfaction) donnent du sens aux
coûts ; d'autre part, sans traçabilité — précisément le vide que comble GIA au niveau branche — il
est difficile d'argumenter chiffres à l'appui.

\section{Lecture pour le projet GIA}
Ces exemples montrent que le MEF et les outils opérationnels (ticketing, tableaux de bord) sont
\textbf{complémentaires} : le premier cadre un diagnostic multi-pôles ; le second fournit des
données factuelles sur l'activité de support. L'extension naturelle du présent travail consisterait
à définir, avec les parties prenantes Naftal, un petit nombre d'indicateurs (délais moyens, taux de
tickets résolus, charge par technicien) et à les suivre dans le temps — en cohérence avec les
recommandations issues des évaluations de la fonction SI dans le groupe~\cite{meghari2014}.
